\subsection{مبانی نظری و فرمول‌بندی ریاضی}
فرض کنید ماتریس وزن $\mathbf{W} \in \mathbb{R}^{d_{\text{out}} \times d_{\text{in}}}$ داده شده باشد. ابعاد ورودی و خروجی ابتدا به $d$ عامل صحیح افراز می‌شوند:
\begin{equation}
d_{\text{in}} = \prod_{k=1}^d p_k, \quad d_{\text{out}} = \prod_{k=1}^d q_k
\end{equation}
ماتریس $\mathbf{W}$ به صورت یک تانسور مرتبه $2d$ با ابعاد $(q_1, \dots, q_d, p_1, \dots, p_d)$ بازآرایی و جایگشت می‌یابد تا جفت‌های ورودی-خروجی در کنار یکدیگر قرار گیرند. سپس این تانسور به فرم قطار تانسوری تجزیه می‌شود:
\begin{equation}
\mathbf{W}(i_1, \dots, i_d, j_1, \dots, j_d) = \mathcal{G}_1[i_1, j_1] \cdot \mathcal{G}_2[i_2, j_2] \cdots \mathcal{G}_d[i_d, j_d]
\end{equation}
که در آن $\mathcal{G}_k \in \mathbb{R}^{r_{k-1} \times q_k \times p_k \times r_k}$ هسته‌های تانسوری هستند و رتبه‌های مرزی برابر با $r_0 = r_d = 1$ می‌باشند. مقادیر $\{r_k\}$ به عنوان رتبه‌های قطار تانسوری (\lr{TT-ranks}) شناخته می‌شوند.

تعداد کل پارامترهای ذخیره‌شده در هسته‌های تانسوری برابر است با:
\begin{equation}
N_{\text{TT}} = \sum_{k=1}^d r_{k-1} \cdot p_k \cdot q_k \cdot r_k
\end{equation}
برای دستیابی به کاهش واقعی حجم پارامترها نسبت به ماتریس اولیه ($d_{\text{in}} \cdot d_{\text{out}}$)، رتبه تانسوری باید اکیداً از حد بحرانی کم‌تر باشد:
\begin{equation}
r_{\text{max}} < \frac{d_{\text{out}} \cdot d_{\text{in}}}{\sum_{k=1}^d p_k \cdot q_k}
\end{equation}
در حین گذر پیش‌رو، حاصل‌ضرب تانسور ورودی $\mathbf{X}$ با هسته‌های تانسوری با انقباض متوالی تانسورها (\lr{Tensor Contraction}) با استفاده از دستورات \lr{torch.einsum} یا تبدیل ضرب‌های ماتریسی پی‌درپی محاسبه می‌گردد.
